\chapter{系統分析與設計}
\label{chapter:methodology}

\section{整合運算框架}
\label{sec:framework}

既有的都市微氣候生成設計工作流面臨一個結構性矛盾：高擬真度模擬引擎能提供物理可信的熱壓力場，卻難以在設計迭代所需的時間內提供立即性回饋；而純資料驅動的機器學習代理模型雖可即時推論，卻常缺乏可追溯的地理資訊基礎與物理一致性保證。本研究之整合運算框架回應此一矛盾，由五個依序銜接的核心模組組成，形成一個從地理資訊建構、都市形態模擬、異質圖轉譯、物理資訊深度學習預測，到多目標最佳化與熱舒適人行動線評估的完整閉環流程：

\begin{enumerate}
    \item \textbf{模組一：基本 GIS 資訊建構（真實 OpenStreetMap 建物足跡）}（GIS Foundation Construction）\\
    以環境部智慧城鄉 IoT 微氣象感測網路建立研究場域之座標參考系統與感測器空間分布；並整合 OpenStreetMap 全域建物足跡與 ETH 全球樹冠高度資料（第~\ref{subsec:v4_real_dataset}節），取代早期單一依賴 NLSC 圖資來源之作法，支援 300 個真實場域之大規模自動化篩選與擷取，作為後續模擬與驗證之地理錨點。
    \item \textbf{模組二：都市形態模擬}（Urban Morphology Simulation）\\
    以 Honeybee／Radiance／EnergyPlus 於 Rhino Grasshopper 環境中執行全場域微氣候逐時模擬，模擬對象為真實擷取之都市量體幾何輸入，取得作為機器學習訓練監督信號之熱壓力目標值（V5：官方光線追蹤與熱傳求解引擎，見第~\ref{sec:ladybug_workflow}節）。
    \item \textbf{模組三：運算框架 —— 真實都市建物足跡與樓高擷取}（Real-World Building Footprint \& Height Extraction）\\
    自 NLSC 多重圖資服務融合擷取單一場域之高精度真實建物輪廓與樓層資訊，取代傳統仰賴蒙地卡羅隨機採樣之單一隨機幾何生成方式，與模組一之 OSM 大規模擷取互補，共同使模型驗證得以錨定於實際都市紋理。
    \item \textbf{模組四：異質圖建構與 PI-ST-GNN 推理}（Heterogeneous Graph Construction \& Inference）\\
    將場景幾何轉譯為包含建物節點與空氣評估節點之異質圖，並以物理資訊時空圖神經網路——Physics-Informed Spatio-Temporal Graph Neural Network（PI-ST-GNN，全名於此首次完整標示，後續章節均沿用此縮寫）——低延遲預測全場域 UTCI 熱壓力場。
    \item \textbf{模組五：Grasshopper 介面整合與多目標最佳化}（Multi-Objective Optimization \& Thermal Comfort Walkway Design）\\
    以訓練完成的 GNN 作為代理模型嵌入 NSGA-II 演算法，同步評估都市形態之熱舒適人行動線品質，透過 FastAPI／WebSocket 後端與 Rhino Grasshopper 介面進行即時通訊，支援互動式閉環設計決策。
\end{enumerate}

圖~\ref{fig:module_flow} 以模組方塊與資訊流箭頭呈現上述五模組之整合架構：模組一（真實 OSM 建物足跡＋IoT）與模組三（NLSC 高精度單場域擷取）為兩條互補之地理資訊輸入路徑，共同匯入模組二之微氣候模擬；模組二輸出之熱壓力標籤驅動模組四之異質圖建構與 PI-ST-GNN 訓練／推理；模組四訓練完成之代理模型最終嵌入模組五，於 Grasshopper 介面內支援 NSGA-II 多目標最佳化與熱舒適人行動線設計之即時回饋。

\begin{figure}[htbp]
    \centering
    \begin{tikzpicture}[
        node distance=8mm and 10mm,
        box/.style={rectangle, rounded corners=2pt, draw=black, thin,
            minimum width=3.05cm, minimum height=1.35cm, align=center,
            font=\small, text=black, fill=white},
        m1/.style={box},
        m2/.style={box},
        m3/.style={box},
        m4/.style={box},
        m5/.style={box},
        lbl/.style={font=\scriptsize\itshape, text=black!55},
        arr/.style={-{Stealth[length=2.6mm]}, thick, draw=black!70},
        darr/.style={-{Stealth[length=2.6mm]}, thick, dashed, draw=black!55}
    ]
        % row 1: two complementary GIS inputs
        \node[m1] (m1) {模組一\\基本 GIS 資訊建構\\{\scriptsize 真實 OSM＋IoT}};
        \node[m3, right=of m1] (m3) {模組三\\真實建物擷取\\{\scriptsize NLSC 三服務融合}};

        % row 2: simulation, converges from both
        \node[m2] (m2) at ($(m1)!0.5!(m3) + (0,-2.1)$) {模組二\\都市形態模擬\\{\scriptsize Honeybee／Radiance／EnergyPlus}};

        % row 3: graph + model
        \node[m4, below=14mm of m2] (m4) {模組四\\異質圖建構與\\PI-ST-GNN 推理};

        % row 4: optimization / design
        \node[m5, below=14mm of m4] (m5) {模組五\\Grasshopper 介面整合\\多目標最佳化};

        \draw[arr] (m1) -- (m2) node[midway, left=1mm, lbl] {座標／感測校準};
        \draw[arr] (m3) -- (m2) node[midway, right=1mm, lbl] {建物幾何};
        \draw[arr] (m2) -- (m4) node[midway, right=1mm, lbl] {UTCI 監督標籤};
        \draw[arr] (m4) -- (m5) node[midway, right=1mm, lbl] {訓練完成之代理模型};
        \draw[darr] (m5.west) .. controls +(-2.6,0) and +(-2.6,0) .. (m1.west)
            node[midway, left=1mm, lbl, align=center] {設計迭代\\回饋閉環};
    \end{tikzpicture}
    \caption[整合運算框架五模組架構與資訊流]{整合運算框架五模組架構與資訊流：模組一、三為互補之地理資訊輸入路徑，模組二至五依序構成模擬、建模、最佳化之閉環流程（虛線為設計迭代回饋路徑）}
    \label{fig:module_flow}
\end{figure}

圖~\ref{fig:data_pipeline_overview} 進一步以資料架構圖展示上述五個模組間之完整資料流向，並標示各處理階段與本章對應小節之關係。

\begin{figure}[htbp]
    \centering
    \includegraphics[width=\textwidth]{../urban-thermal-gnn/04_training/figures/fig_data_pipeline_overview.png}
    \caption[整合運算框架資料架構圖]{整合運算框架資料架構圖}
    \label{fig:data_pipeline_overview}
\end{figure}

上圖由左上四項資料來源（GIS 建物擷取、IoT 感測器、CWA 氣象站、隨機生成幾何）出發，經都市形態模擬、異質圖建構，匯入 PI-ST-GNN 推理引擎，最終驅動下方三項設計決策輸出（NSGA-II 最佳化、熱舒適人行動線設計、Grasshopper 介面回饋），色塊對應本章各小節。

\subsection{設計動機與限制條件}
\label{subsec:framework_motivation}

上述五模組之劃分並非任意的工程切割，而是直接對應第~\ref{chapter:intro}章第~\ref{sec:problem}節所揭示之三項研究缺口與其各自之限制條件。模組一、二、三共同回應「資料是否錨定於真實都市紋理」此一限制條件——若訓練資料仍停留於程序化隨機生成之封閉解空間，模型即便在自身測試集上取得高決定係數，仍難以回答「這個模型是否適用於真實案址」之根本問題，此即為何本研究於真實場景資料集（第~\ref{sec:v4_validation}節）中將建物幾何、樹冠、氣象輸入與微氣候模擬引擎全面替換為 OpenStreetMap、ETH 全球樹冠高度、環境部 IoT 逐時實測資料，以及 Radiance／EnergyPlus 光線追蹤與熱傳求解（V5），而非僅止於程序化模擬或近似式引擎的自我驗證。模組四回應「單次推理延遲是否落在設計迭代可承受之量級」此一限制條件，這是決定代理模型能否被實際嵌入 NSGA-II 演算法內迴圈、而非僅作為離線批次分析工具的關鍵前提；若代理模型之單次推理仍需以分鐘計，則第~\ref{sec:nsga2}節所述之數十代演化搜尋將重新落入與高擬真度模擬引擎相同的運算瓶頸，代理模型之存在意義將大打折扣。模組五回應「運算結果是否能實際回饋至建築設計決策」此一限制條件——一個精度再高的預測模型，若缺乏與 Rhino/Grasshopper 等建築師慣用設計環境的介面整合，其產出便難以直接進入方案發展流程；此一限制條件正是本研究選擇以 FastAPI／WebSocket 即時通訊、而非離線批次匯出報表作為系統輸出形式的核心理由。

\subsection{模組化之必要性論證}
\label{subsec:framework_modularity}

本研究刻意採用模組化架構，而非將地理資訊擷取、模擬取樣與模型訓練整合為單一不可拆解的端到端流程，原因有三。其一，\textbf{資料來源之可替換性}：模組一至三所處理之地理資訊來源本質上具有高度異質性與時效性——政府開放圖資之圖層版本會定期更新、OSM 標籤完整度因地區而異、IoT 測站布建範圍逐年擴增——若將資料擷取邏輯與下游模型訓練緊密耦合，任何資料來源之變動都將牽動整個系統的重新設計；模組化設計使模組三之建物擷取方法（OSM 批次擷取或三服務融合向量化）能被獨立替換或並行使用，而不影響模組四、五之圖神經網路架構與最佳化邏輯。其二，\textbf{驗證責任之可分離性}：模組二（都市形態模擬）之輸出品質決定了模組四訓練標籤之上限，若二者未被清楚劃分為獨立模組，模型預測誤差將難以歸因於「模擬引擎本身之物理近似限制」或「圖神經網路之學習能力不足」，此一可歸因性正是第~\ref{sec:limitations}節局限性討論得以逐條追溯原因、而非籠統歸咎於「模型精度不足」的方法論基礎。其三，\textbf{部署環境之可攜性}：模組五之 FastAPI／WebSocket 後端與 Grasshopper 前端刻意與模組一至四之訓練流程解耦，使得同一套已訓練完成之代理模型檔案，能在不同硬體環境（訓練用 GPU 工作站與設計端筆記型電腦）之間直接搬遷部署，訓練與推理兩端無需共用相同的資料前處理程式碼路徑。

\subsection{PI-ST-GNN 全名與可解釋性定位}
\label{subsec:framework_interpretability}

PI-ST-GNN（Physics-Informed Spatio-Temporal Graph Neural Network，物理資訊時空圖神經網路）此一命名之三個組成部分，分別對應模型設計中三項刻意保留的可解釋性錨點，而非僅為技術堆疊之標籤。「Physics-Informed」對應第~\ref{sec:physics_loss}節之三項物理資訊軟約束損失（輻射一致性、時序平滑性、風場阻滯），其設計動機並非單純追求數值精度之提升，而是確保模型輸出可被人類專家以物理常識直接檢核——例如陰影區之預測溫度不應高於日照區，此一可檢核性使建築師與口試委員得以在不理解模型內部權重的前提下，仍能對模型輸出之合理性進行判斷。「Spatio-Temporal」對應模型同時處理空間關係型圖卷積（第~\ref{subsec:rgcn}節 RGCN）與時序演化（第~\ref{subsec:lstm}節 LSTM 暖啟動機制）雙重結構，使模型輸出之全日 11 時段熱壓力場彼此之間具有物理上合理的連續性，而非各時步獨立、可能出現突兀跳動的離散預測。「Graph Neural Network」則對應第~\ref{sec:graph_construction}節以異質圖顯式編碼建物與空氣節點間之遮蔽、對流、蒸散、語義、毗鄰五類因果關係邊，使模型之訊息傳遞路徑本身即具備可追溯性——第~\ref{subsec:trustworthiness}節之 GNNExplainer 分析，即是沿著這些顯式邊類型回溯特定預測值之主要貢獻來源，此一可解釋性分析路徑，在純像素卷積或全連接黑盒模型中並不存在對應的結構基礎。

\subsection{與 UrbanGraph 之關係}
\label{subsec:framework_urbangraph_relation}

本研究之 RGCN 關係型聚合公式與 LSTM 暖啟動機制，係採用並延伸自 \citet{xin2025urbangraph} 提出之 UrbanGraph 架構（詳見第~\ref{subsec:rgcn}、\ref{subsec:lstm}節之公式推導），其異質圖之因果拓撲建構邏輯（第~\ref{sec:graph_construction}節之遮蔽邊、對流邊建構規則）亦直接沿用 UrbanGraph 之硬性因果剪枝設計。本研究在此一既有架構基礎上之延伸具體有三：其一，將原架構單一節點類型之設計顯式區分為建物與空氣兩種異質節點類型，並以獨立編碼器分別處理其靜態幾何與動態氣象特徵；其二，於原有圖拓撲硬約束之外，另行引入三項物理資訊軟約束損失作為訓練過程之補充監督機制；其三，將原本可能侷限於單一時步預測之架構，擴展為一次輸出全日 11 時段 UTCI 場之多步驟時序預測，並整合至第~\ref{chapter:results}章之端到端低延遲推理與設計迴圈。上述沿用與延伸之分界，於第~\ref{sec:contributions}節研究貢獻中將以逐項對照方式再次明確標示，避免讀者將既有架構之既有能力誤認為本研究之原創貢獻。

\subsection{設計原則小結}
\label{subsec:framework_principles}

綜合上述四小節，本研究整合運算框架之設計可歸納為三項一貫原則：（一）\textbf{真實性優先於規模}——寧可縮小訓練場景數量（V5 之 300 個真實場景），也優先確保幾何、氣象與物理模擬引擎均錨定於真實世界，而非以程序化生成手段換取更大但脫離真實紋理的資料規模；（二）\textbf{延遲預算驅動架構選型}——模組四之圖神經網路架構選型（RGCN 而非全連接網路、代理模型而非直接呼叫模擬引擎）之根本考量始終是能否落在模組五 NSGA-II 演化搜尋所能負擔之延遲預算內，而非單純追求排行榜式的精度極大化；（三）\textbf{可解釋性作為架構約束、而非事後補充}——第~\ref{subsec:framework_interpretability}節所述之物理資訊損失與顯式異質邊設計，於架構設計階段即已納入考量，使第~\ref{subsec:trustworthiness}節之可解釋性分析得以有明確的結構基礎可供回溯，而非在模型訓練完成後才回頭嘗試以事後方法（post-hoc）解釋一個純黑盒模型的行為。


\section{模組一：基本 GIS 資訊建構}
\label{sec:gis_foundation}

本研究以政府開放圖資作為建構研究場域地理資訊基礎之資料來源，在不依賴商用測繪資料授權的前提下，取得具備樓層屬性的建物量體資料，並建立一套可重現的 GIS 資訊擷取流程。圖~\ref{fig:site_context} 由國家尺度逐層放大至研究場域尺度，說明研究場域在真實地理脈絡中的定位。

\begin{figure}[htbp]
    \centering
    \includegraphics[width=\textwidth]{../urban-thermal-gnn/04_training/figures/fig_geo1_site_context.png}
    \caption[研究場域地理脈絡]{研究場域地理脈絡：由國家尺度逐層放大至場域尺度}
    \label{fig:site_context}
\end{figure}

上圖 (a) 為台灣全島尺度，紅框標示真實場景資料集之擴大研究區域、紅星標示陽明交大光復校區錨點；(b) 為擴大研究區域之真實 IoT 感測器（藍點，2,018 站）與 300 個選定真實訓練場景（綠點）之空間分布；(c) 為一個真實 $80\times80$ 公尺場景之幾何：真實 OSM 建物量體、ETH 樹冠與空氣評估節點格點——取代早期版本之程序化示意場景，直接以真實地理幾何呈現場域尺度。

\subsection{建物圖資取用流程與場景篩選演算法}
\label{subsec:crs_sources}

場景幾何統一以 WGS84 地理座標（EPSG:4326）為基準，並於各場景局部運算時轉換為以場景中心為原點之公尺制平面座標。訓練資料所需之建物圖資面臨一核心約束：須在數百個場景之區域尺度上批次、可重現地取得真實建物足跡與樓高，而互動式測繪服務（如國土測繪中心 WMTS／WMS，須逐圖磚人工套疊判讀）難以支撐此一取樣規模。為此，本研究之真實場景資料集改以 OpenStreetMap（OSM）為建物與地表圖資來源，其離線區域擷取機制與樓高判讀邏輯詳見第~\ref{sec:building_extraction}節；本節聚焦於場景候選之篩選演算法。

場景候選之產生並非於區域內任意佈點，而是以真實環境部 IoT 測站之座標為中心，使每一候選場景中心恰為一真實測站；候選經最小間距 150 公尺之貪婪去重，確保 $80\times80$ 公尺場景互不重疊。篩選演算法對每一候選逐一檢核四項條件（真實建物、地表異質性、真實 IoT 站點、真實樹冠，詳見第~\ref{subsec:v4_real_dataset}節），全數通過者再依綜合品質評分（真實樓高建物數、IoT 密度、樹冠像元數、建物密度加權）排序，取前 300 名構成訓練場景集。此一「以真實測站為中心、以四項條件過濾、以品質評分排序」之演算法，使取樣同時滿足感測器校準之近接性與真實幾何之完整性，其結果之空間分布見第~\ref{subsec:v4_real_dataset}節。


\subsection{IoT 微氣象感測器空間資料}
\label{subsec:iot_spatial}

除靜態建物圖資外，本研究整合環境部智慧城鄉空品微型感測器網路之即時觀測資料，作為模擬邊界條件校準、氣象驅動與研究場域選定之依據。為使真實場景資料集（第~\ref{subsec:v4_real_dataset}節）達科學可信之取樣數，感測器資料涵蓋範圍由早期版本之新竹市 15 公里圈擴大至以陽明交大光復校區為錨點、涵蓋新竹市／新竹縣／桃園市／苗栗縣／臺中市之研究區域，區域內共納入 2,018 站真實感測器，其空間分布與逐時量測特性如圖~\ref{fig:iot_sensor_map} 所示。

\begin{figure}[htbp]
    \centering
    \includegraphics[width=\textwidth]{../urban-thermal-gnn/04_training/figures/fig_iot_sensor_map.png}
    \caption[IoT 微型感測器空間分布實測圖]{真實場景資料集擴大研究區域之環境部 IoT 微型感測器空間分布實測圖（區域內逐時觀測，2025 年夏季正午時段）}
    \label{fig:iot_sensor_map}
\end{figure}

上圖 (a) 為空氣溫度分布，色階由綠（低溫）至紅（高溫），可見都市化核心與丘陵地帶顯著溫差；(b) 為相對濕度分布，沿海與內陸呈現梯度差異；紅色星號標示研究區域錨點（陽明交大光復校區）。

在感測器密度分析階段，本研究以感測器站點之空間座標本身為對象，並以核密度估計（Kernel Density Estimation, KDE）評估其空間分布密度，用以界定真實場景之候選取樣區域。此一感測器密度分析與最終選定之 300 個真實場景之對照，如圖~\ref{fig:studyarea} 所示。

\begin{figure}[htbp]
    \centering
    \includegraphics[width=0.95\textwidth]{../urban-thermal-gnn/04_training/figures/fig_studyarea_sensor_density.png}
    \caption[研究範圍選定]{真實場景資料集研究範圍選定：擴大區域之 IoT 感測器密度曲面與 300 個選定真實場景對照}
    \label{fig:studyarea}
\end{figure}

上圖 (a) 為擴大研究區域內之核密度估計曲面，青點為 2,018 站真實感測器、綠點為 300 個選定之真實訓練場景、紅星為錨點；(b) 為錨點核心區之局部放大，呈現選定場景與周邊真實感測器之空間對應關係。


\subsection{全球樹冠高度資料}
\label{subsec:canopy_data}

建物與感測器圖資之外，喬木樹冠高度是異質圖中 \texttt{veg\_et}（植被蒸散）邊緣與空氣節點 $T_h$（最近樹木高度影響）特徵的真實資料來源。本研究採用 ETH Zurich 發布之全球樹冠高度資料集 ETH GlobalCanopyHeight（10 公尺解析度、2020 年版本，WGS84 地理座標）\citep{lang2023high}，該資料集以 Sentinel-2 多光譜影像結合 GEDI 星載光達（Spaceborne LiDAR）訓練深度學習迴歸模型反演逐像元樹冠高度，覆蓋全球陸地表面。研究場域所在圖磚（N24E120）之空間範圍與研究場域錨點之對照如圖~\ref{fig:canopy_map} 所示。

\begin{figure}[htbp]
    \centering
    \includegraphics[width=\textwidth]{../urban-thermal-gnn/04_training/figures/fig_canopy_height_map.png}
    \caption[樹冠高度資料與研究場域對照]{ETH GlobalCanopyHeight 樹冠高度資料與研究場域對照}
    \label{fig:canopy_map}
\end{figure}

上圖 (a) 為研究場域周邊 15 公里範圍內之樹冠高度分布，疊繪於真實衛星影像底圖之上，色階採與 ETH 官方線上展示工具一致之 inferno 色彩尺度（深紫至亮橙代表樹冠高度遞增），虛線圓圈與圖~\ref{fig:iot_sensor_map} 採相同之 15 公里範圍，供跨圖對照；灰階（無色階疊加）區域代表該處無樹冠訊號（NoData），對應道路、建物等不透水鋪面。可見新竹沿海丘陵地帶具有大面積連續天然林覆蓋（樹冠高度可達 30 公尺以上），與都市化核心區之低樹冠覆蓋形成鮮明對比。(b) 為場域局部放大，青框標示 $80\times80$ 公尺研究場域之實際地理位置與尺度；可見研究場域本身座落於高密度建成核心，該 $80\times80$ 公尺範圍內之像元均為無樹冠訊號，但周邊數十至數百公尺外仍可見明顯樹冠覆蓋。

在資料整合流程上，本研究依循 sensing\_integration/canopy\_loader.py 之取樣邏輯：對每一棵樹木節點之場域局部座標 $(x, y)$，先轉換為 WGS84 經緯度，再以最近像元法（Nearest-neighbor Sampling）自 GeoTIFF 讀取對應樹冠高度值；若該位置存在有效觀測值（高度 $\ge 1$ 公尺），則以此真實觀測值覆蓋原先由隨機幾何採樣器生成之樹木高度與推算樹冠半徑，並標記資料來源為 \texttt{canopy\_map}；若該位置無有效觀測，則保留隨機生成之高度值，標記為 \texttt{random\_fallback}。此一「真實資料優先、程序化生成為輔」的分層取樣策略，使異質圖之樹木節點特徵在有真實觀測依據之處錨定於實測數據。


\subsection{真實場景訓練資料集}
\label{subsec:v4_real_dataset}

為使模型訓練直接立基於真實都市地理資訊而非全然程序化生成之幾何，本研究於前述 GIS 圖資基礎上，建立一套完全以真實資料驅動之 300 場景訓練資料集。此一真實場景之選址、地理資料與氣象邊界條件，同時是本研究最終採用之 V5 主要模型與其前身（V1--V4，見第~\ref{sec:v4_validation}節）共用之基礎——僅下游之微氣候模擬引擎（第~\ref{sec:ladybug_workflow}節）與異質圖動態邊建構方法（第~\ref{subsec:edges}節）有所差異，場址、建物、樹冠與 IoT 氣象等地理資料本身並無版本間差異。其選址以陽明交通大學光復校區為錨點，為達科學可信之取樣數而將搜尋範圍擴及新竹市、新竹縣、桃園市、苗栗縣與臺中市；候選場址直接取自環境部 IoT 測站之真實座標（去重至最小間距 150 公尺，避免 $80\times80$ 公尺場域重疊），並逐一檢核以下四項條件，僅保留同時滿足全部條件者：

\begin{enumerate}
    \item \textbf{真實建物幾何（陰影訊息）：}場域 120 公尺陰影脈絡範圍內至少含 3 棟真實 OSM 建物足跡，且至少 1 棟位於場域中心 70 公尺內，確保場景具真實建物量體與陰影結構；
    \item \textbf{地表覆蓋異質性：}場域周邊具建成（built）以外之至少一類地面材質（草地、水體、裸土、鋪面等，取自 OSM \texttt{landuse}／\texttt{natural}／\texttt{leisure}），合計 $\ge 2$ 類物理表面，提供反射率與蓄熱異質性；
    \item \textbf{高密度真實 IoT：}場域中心即為一真實環境部微型感測器站點（最近測站距離 0 公尺），可直接以該站真實觀測驅動與校準；
    \item \textbf{真實植栽陰影：}場域 $80\times80$ 公尺範圍內具 ETH GlobalCanopyHeight 有效樹冠像元，確保含真實植栽遮蔽訊息。
\end{enumerate}

全區共 304 個候選場址通過全部條件，依綜合品質評分（真實樓高建物數、IoT 密度、樹冠像元、建物密度加權）取前 300 名，其於真實地圖上之分布如圖~\ref{fig:selected_sites} 所示。300 場景依序輪派至 6、7、8 三個月份（各 100 場景），使資料集涵蓋整個實測夏季。各項資料真實度說明如下：建物\textbf{足跡}為 100\% 真實（每場景平均 11.3 棟，直接來自 OSM）；建物\textbf{樓高}則因臺灣 OSM 高度標籤稀疏，僅約 14\% 建物具真實 \texttt{height}／\texttt{building:levels} 標籤，其餘以三層（10.8 公尺）之樓層預設值回填，屬本資料集之已知限制；喬木高度取自 ETH 樹冠實測；空氣溫度與相對濕度則為\textbf{全部 300 場景}各自所在 IoT 測站之真實逐時觀測（溫、濕度覆蓋率均為 300/300）。

\begin{figure}[htbp]
    \centering
    \includegraphics[width=\textwidth]{../urban-thermal-gnn/04_training/figures/fig_selected_real_sites.png}
    \caption[真實訓練場景選取結果]{真實訓練場景選取結果：(a) 300 個場景於擴大搜尋範圍（新竹／桃園／苗栗／臺中，錨定陽明交大光復校區）之區域分布，點色對應綜合選址評分；(b) 錨點核心區之場景，$80\times80$ 公尺場域邊界按真實比例繪製於 OpenStreetMap 底圖之上。每一場景同時滿足真實建物、地表異質性、真實 IoT 站點與真實樹冠四項條件。}
    \label{fig:selected_sites}
\end{figure}

為呈現訓練資料所涵蓋之真實都市幾何多樣性，圖~\ref{fig:scene_array} 自 300 個場景中取 25 個清晰易讀者以 $5\times5$ 陣列並列，各場景之建物量體均為真實 OSM 足跡（色階對應真實簷高），綠點為 ETH 樹冠實測植栽；可見場景涵蓋高密度建成、街廓邊緣、疏落郊區等不同都市紋理與棟距分佈。

\begin{figure}[htbp]
    \centering
    \includegraphics[width=0.96\textwidth]{../urban-thermal-gnn/04_training/figures/fig_scene_array_25.png}
    \caption[真實訓練場景代表性樣本]{300 個真實訓練場景中之 25 個代表性場景（$5\times5$ 陣列）：真實 OSM 建物足跡（色階＝簷高）與 ETH 樹冠。各子圖標題示場景編號、建物棟數與所屬月份。}
    \label{fig:scene_array}
\end{figure}

圖~\ref{fig:scene_array} 所示之靜態幾何多樣性，於逐時模擬後進一步轉譯為隨時間演變之熱壓力場。圖~\ref{fig:heatmap_grid} 取其中 9 個場景（依喬木數量遞增排序），以 V5 Radiance／EnergyPlus 模擬引擎（第~\ref{sec:ladybug_workflow}節）呈現 08:00 至 18:00 六個代表時步之 UTCI 熱力圖並列結果，各時步之數值皆為該場景所屬月份之模擬代表值，非跨月平均。可見隨喬木數量增加，場域內低溫（藍綠）區域範圍隨之擴大；各場景之熱壓力峰值時段普遍集中於太陽仰角最高之 10:00--14:00，午後（16:00--18:00）則隨太陽輻射減弱而回落——此一逐時演變之空間熱壓力型態，即為第~\ref{sec:graph_construction}節異質圖動態邊（\texttt{shadow}、\texttt{veg\_et}、\texttt{convective}）所需學習之核心監督訊號來源。

\begin{figure}[htbp]
    \centering
    \includegraphics[width=\textwidth]{../urban-thermal-gnn/04_training/figures/fig_v5_heatmap_grid_9x6.png}
    \caption[V5 模擬熱力圖陣列]{V5（Radiance／EnergyPlus）模擬之 9 場景 $\times$ 6 時步 UTCI 熱力圖陣列：場景由上至下依喬木數量遞增排序，各列標示場景編號與建物／喬木棟數；灰色為真實 OSM 建物足跡，綠色圓為 ETH 樹冠實測植栽；各時步之數值皆為該場景所屬月份之模擬代表值，非跨月平均。}
    \label{fig:heatmap_grid}
\end{figure}


\section{模組二：都市形態模擬——物理式微氣候逐時模擬}
\label{sec:ladybug_workflow}

微氣候逐時模擬以 Ladybug Tools 生態系之 Honeybee／Radiance／EnergyPlus 官方 \texttt{utci\_comfort\_map} recipe 實作，逐一評估場景內每個空氣節點之逐時熱壓力目標值 $T_a, \text{MRT}, V_a, \text{RH}, \text{UTCI}$。其運算核心包含四個相互耦合之真實物理求解模組：

\begin{enumerate}
    \item \textbf{天空可視因子（Sky View Factor, SVF）}：以 Radiance \texttt{sky\_view} recipe 執行光線追蹤，自各節點向天空半球發射大量取樣射線，依建物與喬木之三維量體幾何計算實際被遮擋之立體角比例。
    \item \textbf{短波輻射與陰影}：以 Radiance 二階段（Two-phase）方法求解直達、漫射與地面反射三分量短波輻射通量，取代單純的幾何陰影布林判定，逐時反映太陽位置、大氣條件與周邊量體之真實光線遮蔽關係。
    \item \textbf{戶外平均輻射溫度（MRT）}：短波分量取自上述 Radiance 求解結果，長波分量則以 EnergyPlus 求解建物表面溫度後，透過 \texttt{ladybug-comfort} 之 \texttt{OutdoorSolarCal} 合成，為真實建物外殼熱傳導與蓄熱行為之直接輸出，而非經驗式近似。
    \item \textbf{風場}：現行版本以每一場景單一之參考風速代表全場域風況（見下段殘餘限制說明），尚未逐點求解街谷尺度風場遮蔽。
\end{enumerate}

UTCI 直接取自 \texttt{utci\_comfort\_map} recipe 之原生多項式輸出（不再另行呼叫 \texttt{pythermalcomfort} 套件重新合成）。模擬時段涵蓋 08:00 至 18:00 共 11 個整點時步（$T = 11$），捕捉日間熱壓力最嚴峻之逐時變化軌跡。

需說明者：本研究之微氣候模擬引擎歷經 V1--V4 之自製 Python 物理近似模型（幾何射線近似 SVF 與陰影、地表能量平衡經驗式 MRT，\textbf{未}呼叫任何外部模擬套件），至 V5 起正式改採上述 Radiance 6.2.1（光線追蹤）與 EnergyPlus 25.1.0（建物熱傳求解）之官方模擬流程——此一轉換之完整比較與 V1--V4 近似引擎之瑕疵揭露，見第~\ref{sec:v4_validation}節與附錄~\ref{appx:dynamic_edges_fix}。V5 現行版本仍有以下已知殘餘限制：（一）風場以每場景單一參考風速代表，未逐點求解街谷尺度繞流；（二）建物外殼構造採 ASHRAE 標準構造假設，而非逐棟真實材料實測；（三）陰影判定之來源索引（哪一建物／喬木造成特定節點陰影）為幾何後處理之推導代理值，非 recipe 原生輸出。上述限制均不影響 UTCI 本身之物理求解性質，但為後續研究可精進之方向，詳見第~\ref{sec:limitations}節。

氣象邊界條件於第~\ref{subsec:v4_real_dataset}節之真實場景資料集中，直接採用各場景所在地之環境部 IoT 測站真實逐時氣溫與相對濕度（經合理性驗證，故障測站退回中央氣象署區域觀測）、中央氣象署測站真實風速，以及區域 TMY（Typical Meteorological Year）EPW 之實測日射剖面（因 CWA 該測站未提供日射觀測，故日射採區域典型氣象年之實測輻射，而非場景當日觀測），並由 \texttt{13\_build\_scenario\_epw\_v5.py} 逐場景合成對應之客製化 EPW 檔案供 Radiance／EnergyPlus 讀取；此一設計使模擬之熱壓力目標值錨定於真實觀測氣候。

此一模擬流程對每個場景輸出逐時之全場域 UTCI，即為機器學習之監督標籤。圖~\ref{fig:timestep_utci} 以四個真實場景示範此一熱舒適訓練資料之生成：各列為一場景，各行為 08:00 至 18:00 之代表時步，節點顏色為 Radiance／EnergyPlus 模擬之逐時 UTCI。可清楚觀察到熱壓力於日間之累積與消退（清晨低、正午高、傍晚回落），以及真實建物與植栽在各時步所投射之陰影所形成之持續性空間結構——此即時空圖神經網路所需學習之「幾何—時序—熱壓力」映射關係之資料來源。

\begin{figure}[htbp]
    \centering
    \includegraphics[width=\textwidth]{../urban-thermal-gnn/04_training/figures/fig_v5_timestep_utci.png}
    \caption[熱舒適訓練資料之逐時生成]{熱舒適訓練資料之逐時生成（V5，Radiance／EnergyPlus 模擬）：真實場景於 08:00--18:00 各代表時步之模擬 UTCI 場。展現熱壓力之日間動態與真實幾何遮蔽所致之空間梯度。}
    \label{fig:timestep_utci}
\end{figure}

此一模擬工作流之輸出，除作為機器學習訓練之監督信號外，亦是第~\ref{sec:nsga2}節熱舒適人行動線評估與多目標最佳化之基準比對依據。


\section{模組三：真實都市建物足跡與樓高擷取——開放街圖（OSM）真實建物與地表鋪面圖資擷取}
\label{sec:building_extraction}

\subsection{既有隨機生成方法之侷限}
\label{subsec:parametric_limitation}

模型訓練所需的大規模場景資料集，傳統上由隨機幾何採樣器（Stochastic Geometry Sampler）以蒙地卡羅法在容積率（FAR）、建蔽率（BCR）與縮退規範之法規約束下生成：每個場景包含 2--5 棟建築（矩形或 L 形平面、面寬與進深各 5--25m、樓層數 2--12 層）與 2--5 棵喬木（高度 4--12m），共產生 300 個合法場景。此方法能有效覆蓋法規允許之形態解空間，但其產出之建物量體、棟數關係、街廓紋理與地表覆蓋均為程序化生成，與真實都市紋理之棟距、街廓尺度、樓高分布與鋪面異質性統計特性未必一致，使模型於真實場地應用時的一般化可信度難以驗證。因此，本研究以開放街圖（OpenStreetMap, OSM）之真實建物足跡、樓高與地表覆蓋圖資，全面取代訓練場景之程序化幾何輸入，詳見以下兩小節；既有隨機生成方法則保留作為第~\ref{sec:v4_validation}節之對照資料集。圖~\ref{fig:scenarios} 示例六個典型隨機生成場景，可見其建築量體雖滿足法規約束，但棟距、朝向與棟數關係呈現均質化之程序化特徵，與圖~\ref{fig:scene_array} 之真實場景紋理形成對照。

\begin{figure}[htbp]
    \centering
    \includegraphics[width=\textwidth]{../urban-thermal-gnn/04_training/figures/fig_geo4_scenarios.png}
    \caption[隨機生成場景示意]{六個典型隨機生成場景之建築量體與喬木配置示意}
    \label{fig:scenarios}
\end{figure}


\subsection{OSM 離線區域建物足跡與樓高擷取}
\label{subsec:osm_building_extraction}

以線上 Overpass API 逐場景查詢建物圖資，在 300 個場景之取樣規模下面臨兩項限制：公開服務之流量限制使查詢頻繁逾時或遭拒，且每次查詢之網路延遲使批次取樣耗時不可控。本研究因此改採離線區域擷取：一次性下載 Geofabrik 釋出之臺灣區域 \texttt{.osm.pbf} 封裝檔，以 pyosmium 函式庫解析全部建物與地表覆蓋要素於研究區域邊界框內者，建立含 STRtree 空間索引之常駐記憶體資料結構；此後任一場景之建物與材質查詢即為純記憶體之空間索引查詢，無需再次連線。

建物幾何之萃取僅處理封閉方式（Closed Way）——即端點座標相同、由 OSM 節點座標依序連接而成之多邊形——此涵蓋研究區域內絕大多數建物；由關係（Relation）組成之少數大型建物複合體多邊形未納入處理範圍，如實計入未擷取項目，而非以近似輪廓填補。建物樓高之判讀依序嘗試 \texttt{height}、\texttt{building:height} 標籤，其次以 \texttt{building:levels} 樓層數乘以標準層高（3.6 公尺）推算；當標籤含分號分隔之多值（如合併建物之 \texttt{building:levels="11;10;12"}，反映該建物實際由不同樓高之區塊合併登錄）時，取各值之平均。當上述來源皆缺失時，以三層樓（10.8 公尺）之預設值回填，並以布林旗標標記該筆高度之真實性——僅當高度直接來自 OSM 標籤時計為真實，回填值不計入；此一標記機制使第~\ref{subsec:v4_real_dataset}節得以逐場景統計真實樓高之比例。


\subsection{OSM 地表鋪面與材質圖資擷取}
\label{subsec:osm_material_extraction}

建物幾何之外，地表材質之空間異質性直接決定各節點之反射率（Albedo）與蓄熱特性，是都市微氣候模擬不可或缺的邊界條件。本研究以 OSM 之 \texttt{landuse}、\texttt{natural}、\texttt{leisure} 三類標籤，將地表覆蓋歸類為草地／植被、裸土、鋪面／混凝土、瀝青／工業、水域五種物理材質：\texttt{landuse=grass／meadow／park／forest} 等歸為草地，\texttt{landuse=residential／commercial／retail} 等歸為鋪面，\texttt{landuse=industrial} 歸為瀝青，\texttt{natural=water} 與 \texttt{leisure=swimming\_pool} 等歸為水域，\texttt{landuse=farmland／construction} 等歸為裸土；此一標籤對應表直接沿用第~\ref{subsec:crs_sources}節之 OSM 離線區域索引，與建物擷取共用同一次區域資料載入。

每一材質多邊形依此對應表轉換為地表覆蓋代碼，並以 1 公尺網格柵格化為場景之 \texttt{land\_cover\_map}：地面材質先行填入，建物足跡範圍隨後覆蓋為「建物（built）」代碼，反映屋頂材質優先於其下方地面材質之物理邏輯。此一代碼直接對應第~\ref{sec:ladybug_workflow}節模擬引擎之反射率查表（建物／鋪面 0.30、裸土 0.40、水域 0.00、草地 0.20），取代既有版本中僅以單一場地平均反射率參數近似地表異質性之作法。圖~\ref{fig:material_zones} 以一個真實場景示範此一取用流程：(a) 為 OSM 原始材質多邊形與建物足跡疊圖；(b) 為柵格化後實際餵入模擬引擎之 \texttt{land\_cover\_map}。

\begin{figure}[htbp]
    \centering
    \includegraphics[width=\textwidth]{../urban-thermal-gnn/04_training/figures/fig_material_zones.png}
    \caption[OSM 地表鋪面與材質圖資擷取]{OSM 地表鋪面與材質圖資擷取：(a) 原始材質多邊形與建物足跡；(b) 柵格化 land\_cover\_map。}
    \label{fig:material_zones}
\end{figure}

須說明者：OSM 之 \texttt{landuse} 標註並非對地表覆蓋之窮盡標記，如圖~\ref{fig:material_zones}(b) 所示，場景內相當比例之網格未落於任何材質多邊形範圍內（多對應道路、未標註之鋪面或私有地面）；此類未標註網格於模擬時退回以場域平均反射率參數近似，而非誤植為特定材質，此一取捨保守地避免了對未知地表覆蓋之武斷分類。


\subsection{三服務融合建物向量化方法（單一場地案例驗證用途）}
\label{subsec:vectorize_method}

OSM 離線區域擷取雖能支撐 300 場景規模之批次訓練資料建構，其建物樓高標籤覆蓋率有限；對於需要高精度樓高驗證之單一場地案例研究（第~\ref{chapter:results}章），本研究另建立一套以政府開放圖資為來源之三服務融合建物向量化方法，作為 OSM 之互補驗證來源。真實建物量體之擷取面臨一項核心限制：單一圖資來源難以同時提供精確幾何輪廓與樓層屬性。TOPO01K 圖層雖標註樓高文字，但其柵格化 hatch 填充在高縮放層級下受制圖線寬影響而破碎為大量細小連通元件，直接以連通元件分析取得的建物遮罩並不可靠；B5000 與 EMAP 雖幾何品質較佳，卻不含樓層資訊。本研究因此設計一套三服務交叉驗證方法，讓三個各有所長、各有所缺的圖資來源互補：

\begin{enumerate}
    \item \textbf{建物範圍重建：}將 TOPO01K 之粉紅色填色像素與黑色線繪像素（外框、內部 hatch、樓高文字）聯集為單一「建物範圍遮罩」，而非僅取粉紅色填色本身 —— 因為 hatch 線條與填色像素在製圖時緊密交疊，聯集遮罩可還原被 hatch 切割破碎的填色區域至建物尺度之完整連通元件，且無需透過形態學膨脹等會扭曲邊界的操作。
    \item \textbf{粗細線分離：}以形態學開運算（Morphological Opening）對建物範圍內之黑色線像素進行處理，去除短而細之斜向 hatch 筆劃，保留較粗、較連續之樓高分區邊界線，藉此將單一建物範圍細分為對應各自獨立樓高之子區塊。
    \item \textbf{子區塊輪廓萃取：}對每一子區塊以 OpenCV \citep{bradski2000opencv} 實作之邊界追蹤演算法（cv2.\allowbreak findContours，基於 Suzuki \& Abe 之拓樸邊界追蹤法 \citep{suzuki1985topological}）取得次像素封閉輪廓，再以 Douglas--Peucker 演算法 \citep{douglas1973algorithms} 化簡為平滑多邊形，作為該樓高分區之向量化幾何輸出。
    \item \textbf{多字元群集樓高辨識：}掃描每一子區塊內部之黑色文字像素群集（依尺寸與長寬比篩除細長之殘留 hatch 筆劃），對每一候選群集個別以 Tesseract OCR 引擎 \citep{smith2007overview} 執行光學文字辨識，取代僅於整棟建物幾何中心點單次辨識之作法 —— 大型建物複合體內常包含多個並列之不同樓高分區（如 \texttt{6R}、\texttt{5R}、\texttt{4R} 並列），單次中心點辨識幾乎必然錯過實際文字位置。
    \item \textbf{跨服務信心度評估：}將每一子區塊之向量化輪廓分別與 B5000 黑線封閉區域、EMAP 淡紫灰填色區域進行像素重疊率比對，依兩項驗證來源之通過情形標記高／中／低信心度，供後續篩選使用。
\end{enumerate}

圖~\ref{fig:gis_fusion_pipeline} 完整呈現上述三服務融合方法之系統流程，包含三項圖資之像素對齊機制、TOPO01K 主處理鏈之四道步驟，以及 B5000／EMAP 驗證分支如何匯聚為最終信心度分級輸出。

\begin{figure}[htbp]
    \centering
    \includegraphics[width=\textwidth]{../urban-thermal-gnn/04_training/figures/fig_gis_fusion_pipeline.png}
    \caption[三服務融合建物向量化流程]{三服務融合建物向量化系統流程圖}
    \label{fig:gis_fusion_pipeline}
\end{figure}

上圖橘色為 TOPO01K 主幾何與樓高擷取鏈（建物範圍重建→粗細線分離→子區塊輪廓萃取→多字元群集樓高辨識）；深灰、紫色分別為 B5000、EMAP 之獨立驗證分支；綠色為最終輸出，附帶信心度分級與完整屬性表。

樓層文字之判讀依循新竹地區都市計畫現況圖之標示慣例：\texttt{NR} 表示地上第 N 層、\texttt{T} 表示地面層、\texttt{NB} 表示地下第 N 層、\texttt{NM} 表示夾層。當同一子區塊經多字元群集辨識比對後仍檢出一種以上互斥樓層標籤時，判定為分區尚未完全切割成功之案例，如實標記為待覆核而非任意保留其中一筆結果，以維持資料集之可追溯性。


\subsection{三類幾何來源之分工關係}
\label{subsec:extraction_vs_parametric}

本研究之三類幾何來源並非互斥替代，而依各自之取樣規模與精度優勢分工：OSM 離線區域擷取（第~\ref{subsec:osm_building_extraction}節）支撐 300 場景大規模真實訓練資料集，以真實足跡與地表覆蓋換取模型訓練分佈與真實都市紋理之統計一致性；三服務融合方法（本節）提供單一場地之高精度樓高驗證，作為第~\ref{chapter:results}章案例研究與 NSGA-II 最佳化之現況基準（Baseline）；既有隨機生成場景則保留作為第~\ref{sec:v4_validation}節之對照資料集，用以檢驗真實地理環境訓練相對於程序化幾何訓練之外部效度差異。第~\ref{chapter:results}章之案例研究以三服務融合方法擷取之真實建物足跡作為場地現況輸入，評估模型於未見過真實幾何配置下之一般化表現。


\section{模組四：異質圖建構與 PI-ST-GNN 推理——異質圖建構流程}
\label{sec:graph_construction}

\subsection{節點類型與特徵設計}
\label{subsec:nodes}

都市微氣候系統中，空間元素的物理特性差異顯著——建築物為不透氣剛體，主導陰影投射與風場阻滯；空氣評估點為連續流體場採樣，承載氣象時序資訊。本研究因此採用\textbf{異質圖}（Heterogeneous Graph）結構，將兩類節點分別賦予不同語義維度的特徵向量：

\noindent\textbf{建物節點（Object Nodes，$\mathcal{V}_{\text{obj}}$）}

建物節點共 $N_{\text{obj}}$ 個，對應場景內所有建築量體，每個節點攜帶 7 維靜態幾何特徵向量（見表~\ref{tab:obj_features}）；喬木則不構成獨立之建物節點，其高度資訊改以空氣節點第 7 維特徵（樹木高度正規化，見表~\ref{tab:air_features}）之形式，隨鄰近空氣評估點一併編碼。

\begin{table}[htbp]
    \centering
    \caption{建物節點特徵定義（7維靜態特徵）}
    \label{tab:obj_features}
    \small
    \begin{tabular}{lllc}
        \toprule
        \textbf{特徵名稱} & \textbf{物理意義} & \textbf{正規化方式} & \textbf{範圍} \\
        \midrule
        \texttt{height / 50}        & 建築簷高（m）                  & 除以 50m 上限     & $[0, 1]$ \\
        \texttt{floors / 12}        & 樓層數                        & 除以 12 層上限    & $[0, 1]$ \\
        \texttt{footprint / 2000}   & 底層投影面積（m$^2$）           & 除以 2000m$^2$    & $[0, 1]$ \\
        \texttt{cx / 80}            & 質心 X 座標（m）               & 除以場地寬 80m    & $[0, 1]$ \\
        \texttt{cy / 80}            & 質心 Y 座標（m）               & 除以場地高 80m    & $[0, 1]$ \\
        \texttt{gfa / 20000}        & 總樓地板面積（m$^2$）           & 除以 20000m$^2$   & $[0, 1]$ \\
        \texttt{is\_L\_shape}       & 平面形態旗標（L形=1，矩形=0）  & 二值              & $\{0, 1\}$ \\
        \bottomrule
    \end{tabular}
\end{table}

\noindent\textbf{空氣評估節點（Air Nodes，$\mathcal{V}_{\text{air}}$）}

空氣評估節點共 $N_{\text{air}} = 6{,}241$ 個，佈設於場地 $79 \times 79$ 個 1.0m 間距格點上，每個節點在每個時步攜帶 9 維動態特徵向量（見表~\ref{tab:air_features}）：

\begin{table}[htbp]
    \centering
    \caption{空氣評估節點特徵定義（每時步 9 維動態特徵）}
    \label{tab:air_features}
    \small
    \begin{tabular}{lllc}
        \toprule
        \textbf{特徵名稱} & \textbf{物理意義} & \textbf{正規化方式} & \textbf{範圍} \\
        \midrule
        \texttt{Ta\_norm}   & 空氣溫度（°C）             & z-score 標準化  & $\sim\mathcal{N}(0,1)$ \\
        \texttt{MRT\_norm}  & 平均輻射溫度（°C）         & z-score 標準化  & $\sim\mathcal{N}(0,1)$ \\
        \texttt{Va\_norm}   & 水平風速（m/s）            & z-score 標準化  & $\sim\mathcal{N}(0,1)$ \\
        \texttt{RH\_norm}   & 相對濕度（\%）             & z-score 標準化  & $\sim\mathcal{N}(0,1)$ \\
        \texttt{SVF}        & 天空視角因子              & 原始值 $[0,1]$ & $[0, 1]$ \\
        \texttt{shadow}     & 陰影遮蔽旗標              & 二值            & $\{0, 1\}$ \\
        \texttt{Bh / 50}    & 最近建築高度影響（m）      & 除以 50m 上限   & $[0, 1]$ \\
        \texttt{Th / 12}    & 最近樹木高度影響（m）      & 除以 12m 上限   & $[0, 1]$ \\
        \texttt{Ts\_norm}   & 地表溫度估算值（°C）       & z-score 標準化  & $\sim\mathcal{N}(0,1)$ \\
        \bottomrule
    \end{tabular}
\end{table}

其中地表溫度 $T_s$ 並非直接量測值，而由下式之對流熱平衡估算計算得出：

\begin{equation}
    T_s = T_a + \frac{\text{GHI}}{h_c}
    \label{eq:surface_temp}
\end{equation}

式中 $\text{GHI}$（W/m$^2$）為太陽總輻照度（來自 EPW 逐時資料），$h_c$（W/m$^2$K）為對流換熱係數（本研究設定為 10 W/m$^2$K）。此估算方法在忽略長波輻射再輻射與儲熱效應的一階近似下，能有效代表都市鋪面之熱輻射貢獻 \citep{oke1988street}。

表~\ref{tab:obj_features}、\ref{tab:air_features} 所定義之特徵向量彼此並非孤立的數字欄位，而是共同描述同一個都市街道剖面中的物理幾何關係：太陽仰角決定建物之陰影投射範圍（對應 \texttt{shadow} 旗標）與喬木冠層之樹蔭覆蓋（對應 \texttt{Th} 最近樹木高度影響），兩者共同決定行人所在位置之熱壓力高低。圖~\ref{fig:street_crosssection} 以側視剖面示意此一關係：同一街谷寬度 $W$、建物高度 $H$ 之幾何條件下，曝曬於陰影與樹蔭範圍外之行人位置可能承受 UTCI $> 38^\circ\text{C}$ 之強烈熱壓力，而位於樹蔭遮蔽下之行人位置則可能低於 $32^\circ\text{C}$；空氣節點沿垂直方向多高度取樣（圖右側藍點），使異質圖得以捕捉此一熱壓力隨高度與遮蔽狀態變化之空間梯度，而非僅以單一地面高度之純量值代表整條街道剖面之熱環境。

\begin{figure}[htbp]
    \centering
    \includegraphics[width=0.92\textwidth]{../urban-thermal-gnn/04_training/figures/fig_street_crosssection.png}
    \caption[都市街道剖面示意：太陽角度、建築陰影、喬木樹蔭與行人熱壓力]{都市街道剖面示意：太陽角度、建築陰影、喬木樹蔭與行人熱壓力之幾何關係。原創繪製，示意性呈現表~\ref{tab:obj_features}、\ref{tab:air_features}所定義特徵向量之物理對應關係，非特定場景之精確幾何。}
    \label{fig:street_crosssection}
\end{figure}


\subsection{邊緣類型與因果拓樸}
\label{subsec:edges}

本研究定義五種關係型邊緣（Relational Edges），各自對應一種明確的物理因果機制（見表~\ref{tab:edge_types}）。與建物節點不同，物件節點（Object Nodes）實際上同時涵蓋建築與喬木兩類實體（第~\ref{subsec:nodes}節），\texttt{shadow} 與 \texttt{veg\_et} 兩者之來源端因此分別限定為建築類與喬木類物件節點，而非泛稱之「建物」：

\begin{table}[htbp]
    \centering
    \caption{異質圖五類關係型邊緣定義}
    \label{tab:edge_types}
    \small
    \begin{tabularx}{\textwidth}{lllX}
        \toprule
        \textbf{邊緣類型} & \textbf{連接方向} & \textbf{建立方法} & \textbf{物理因果意義} \\
        \midrule
        \texttt{shadow}      & 建築物件 $\to$ 空氣 & 幾何陰影投射計算 & 建築遮擋導致輻射熱降低 \\
        \texttt{veg\_et}     & 喬木物件 $\to$ 空氣 & 樹冠覆蓋半徑判斷 & 植被蒸散作用形成冷卻效應 \\
        \texttt{convective}  & 空氣 $\to$ 空氣 & 對流方向性權重   & 風場驅動的對流熱交換 \\
        \texttt{semantic}    & 物件 $\leftrightarrow$ 物件 & 全連接（完整關係） & 建物與喬木之語義形態相互關係 \\
        \texttt{contiguity}  & 空氣 $\leftrightarrow$ 空氣 & KNN（$k = 8$）  & 空間毗鄰性熱梯度傳導 \\
        \bottomrule
    \end{tabularx}
\end{table}

其中 \texttt{semantic} 與 \texttt{contiguity} 為靜態邊，於資料載入階段一次建立並全程共用；\texttt{shadow}、\texttt{veg\_et}、\texttt{convective} 三者則隨太陽方位角、風向逐時步變化，於每個時步重新計算，構成動態邊集合 \texttt{dynamic\_edges[t]}——此三類動態邊之來源索引（哪一個建築造成該空氣節點陰影、哪一棵喬木覆蓋該空氣節點）由幾何後處理腳本重建取得：程序化主線資料集使用 \texttt{11\_build\_dynamic\_edge\_cache.py}，真實場景資料集（V5）則使用 \texttt{16\_build\_dynamic\_edge\_cache\_v5.py}（採真實逐場景太陽方位角與真實逐月中央氣象署風向），完整修正過程記錄於附錄~\ref{appx:dynamic_edges_fix}。圖~\ref{fig:graph_construction} 以真實資料集之場景 0（$N_{\text{obj}}=3$，$N_{\text{air}}=6{,}241$）具體呈現此一異質圖之節點分布、真實建立之 KNN 毗鄰邊局部結構，以及完整節點／邊型別綱要。

\begin{figure}[htbp]
    \centering
    \includegraphics[width=\textwidth]{../urban-thermal-gnn/04_training/figures/fig_graph_construction.png}
    \caption[異質圖建構實例]{異質圖建構實例（取自真實訓練資料場景 0）}
    \label{fig:graph_construction}
\end{figure}

上圖 (a) 為建物節點（紅色量體）與空氣評估節點（依時均 UTCI 著色）之空間分布；(b) 為局部 $18\times18$ 公尺窗口內 KNN($k$=8) 毗鄰邊之真實連接結構，以及該範圍內建物節點向場外其他建物延伸之全連接語義邊；(c) 為異質圖節點／邊型別綱要，對應 9 維空氣節點特徵（含地表溫度）與五類關係型邊緣定義。

圖~\ref{fig:graph_topology} 進一步以另一真實場景（場景 10）之全域視角，呈現 \texttt{semantic} 邊如何在場域內所有建物節點間建立全連接拓樸——不同於圖~\ref{fig:graph_construction}(b) 僅呈現單一建物之局部窗口，此處可清楚看見 $N_{\text{obj}}=4$ 個建物節點兩兩相連之完整語義邊網絡，與其下方 $6{,}241$ 個空氣節點所構成之密集 KNN 毗鄰網格形成尺度對比。圖中另以綠色圓標示場景內之喬木冠層位置作為空間對照，惟依第~\ref{subsec:nodes}節之節點定義，喬木高度僅作為空氣節點特徵（樹木高度正規化）之來源，本身不構成獨立圖節點，故未參與語義邊之連接。

\begin{figure}[htbp]
    \centering
    \includegraphics[width=\textwidth]{../urban-thermal-gnn/04_training/figures/fig_geo5_graph_topology.png}
    \caption[異質圖拓樸全域視角]{異質圖拓樸全域視角（真實訓練資料場景 10）}
    \label{fig:graph_topology}
\end{figure}

上圖 (a) 為完整 $80\times80$ 公尺場域，橘點為 6,241 個空氣節點（KNN 毗鄰邊為求可讀性僅隨機顯示 3,000／49,928 條），橘色量體為建物節點、綠色圓為喬木冠層之空間位置（僅供視覺對照，非圖節點），紫色虛線為建物節點間之全連接語義邊；(b) 為場域左下 $40\times40$ 公尺象限放大，標示邊緣類型與 9 維空氣節點特徵之完整定義。

作為對照，圖~\ref{fig:real_scene_graph} 以第~\ref{subsec:v4_real_dataset}節真實場景資料集中之一個實際場景，展示完全以真實 GIS 圖資建構之異質圖：(a) 之建物量體均為 OSM 真實足跡（色階為真實簷高，實線外框標示具真實樓高標籤者），綠點為 ETH 樹冠實測植栽；(b) 之每個空氣節點依 4 公尺格點佈設，節點顏色為以該場景所在真實 IoT 測站逐時氣溫驅動之模擬尖峰 UTCI，可見建物遮蔽與植栽降溫在空間上形成之熱舒適梯度，即為 PI-ST-GNN 所學習之真實空間結構。

\begin{figure}[htbp]
    \centering
    \includegraphics[width=\textwidth]{../urban-thermal-gnn/04_training/figures/fig_real_scene_graph.png}
    \caption[真實 GIS 圖資建構之異質圖範例]{以真實 GIS 圖資建構之異質圖範例（真實場景資料集）：(a) 真實 OSM 建物與 ETH 樹冠幾何；(b) 空氣節點感測網格與 k-NN 鄰接邊，節點色為真實 IoT 氣溫驅動之模擬尖峰 UTCI。}
    \label{fig:real_scene_graph}
\end{figure}


\subsection{圖索引對齊機制（$+N_{\text{obj}}$ 偏移量）}
\label{subsec:index_offset}

異質圖在送入 RGCN 之前，本研究採用\textbf{節點串接}（Node Concatenation）策略，將建物節點特徵向量 $\mathbf{H}_{\text{obj}} \in \mathbb{R}^{N_{\text{obj}} \times d_{\text{obj}}}$ 與空氣節點特徵 $\mathbf{H}_{\text{air}} \in \mathbb{R}^{N_{\text{air}} \times d_{\text{air}}}$ 沿節點維度串接，形成單一混合節點張量：

\begin{equation}
    \mathbf{H} = \left[ \mathbf{H}_{\text{obj}} ; \mathbf{H}_{\text{air}} \right] \in \mathbb{R}^{(N_{\text{obj}} + N_{\text{air}}) \times d}
    \label{eq:node_concat}
\end{equation}

此機制意味著空氣節點在串接後的全域索引空間中，其索引應整體偏移 $N_{\text{obj}}$。因此，KNN 毗鄰邊（\texttt{contiguity}）在建立時，所有空氣節點的 \texttt{edge\_index} 均需加上偏移量 $N_{\text{obj}}$：

\begin{equation}
    \texttt{edge\_index}_{\text{contiguity}} \mathrel{+}= N_{\text{obj}}
    \label{eq:offset}
\end{equation}

此對齊機制確保跨節點類型之邊緣（如 \texttt{shadow}: 建物 $\to$ 空氣）在統一索引空間中仍能正確尋址目標節點，是整個異質圖能以單一 RGCN 前向傳播處理的關鍵實作細節，完整偏移量計算流程另見附錄~\ref{appx:index_offset_detail}。


\section{物理資訊時空圖神經網路 (PI-ST-GNN) 模型架構}
\label{sec:gnn_model}

\subsection{整體架構概述}
\label{subsec:model_overview}

PI-ST-GNN 模型由以下五個主要子模組組成，依前向傳播順序排列（詳見圖~\ref{fig:arch_pipeline}）：

\begin{enumerate}
    \item \textbf{物件節點編碼器（Object Encoder）：}以兩層 MLP 將 7 維靜態幾何特徵投影至 128 維隱藏表示 $\mathbf{h}_{\text{obj}} \in \mathbb{R}^{N_{\text{obj}} \times 128}$。
    \item \textbf{空氣節點編碼器（Air Encoder）：}以兩層 MLP 將每時步 9 維動態氣象特徵投影至 128 維隱藏表示 $\mathbf{h}_{\text{air},t} \in \mathbb{R}^{N_{\text{air}} \times 128}$，對 $T = 11$ 個時步分別計算。
    \item \textbf{全域情境編碼器（Global Context Encoder）：}將場景整體氣象背景（8 維逐時統計特徵：全場域平均 $T_a, V_a, \text{RH}, \text{GHI}$ 及其差分）以 MLP 編碼為時序情境向量 $\mathbf{c} \in \mathbb{R}^{T \times 192}$，與每時步之空氣節點表示融合，注入全域氣候感知能力。
    \item \textbf{關係圖卷積網路（RGCN）：}在每個時步 $t$，將建物與空氣節點之串接表示送入三層堆疊的 RGCN 塊（RGCNBlock），在五種關係型邊緣上進行異質訊息聚合，輸出包含空間因果資訊的混合節點表示。
    \item \textbf{LSTM 時序建模與輸出頭：}以 Fusion MLP 融合 RGCN 輸出與全域情境向量後，送入單層 LSTM，其初始隱藏狀態 $\mathbf{h}_0$ 由 $t = 0$ 時刻的 RGCN 輸出投影得出，最終由輸出 MLP 預測 11 時步之 UTCI 值 $\hat{\mathbf{u}} \in \mathbb{R}^{N_{\text{air}} \times T}$。
\end{enumerate}

\begin{figure}[htbp]
    \centering
    \includegraphics[width=0.92\textwidth]{../urban-thermal-gnn/04_training/figures/fig_arch_clean.png}
    \caption[整合運算框架示意圖]{PI-ST-GNN 整合運算框架示意圖}
    \label{fig:arch_pipeline}
\end{figure}

上圖左側輸入包含建物幾何特徵（7 維）與多時步空氣節點氣象特徵（9 維 $\times$ 11 步），中部為 RGCN 空間建模與 LSTM 時序融合，右側輸出全場域 UTCI 熱壓力場（11 時段），並以物理資訊損失函數進行約束訓練。

上述五個子模組彼此之銜接關係，以及第~\ref{sec:graph_construction}節之異質圖建構如何作為此模型之前置輸入，完整呈現於圖~\ref{fig:overview_pistgnn}。此圖並將第~\ref{subsec:lstm}節所述「LSTM 初始隱藏狀態由 RGCN 輸出投影得出、而非零向量初始化」之時序展開細節，以下方獨立面板逐時步具體標示——每個時步 $t$ 之異質圖先經 RGCN 完成當時步之空間訊息擴散，其輸出再與前一時步之 LSTM 隱藏狀態 $h_{t-1}$ 共同送入 LSTM 單元更新 $h_t$，如此逐時步展開直至 $t=10$（對應 18:00），最終輸出張量 $(N_{\text{air}}, T{=}11, 256)$ 送入輸出 MLP 完成多時步解碼。此一「空間擴散在前、時序遞迴在後」之逐時步交替結構，是 PI-ST-GNN 得以同時捕捉建物陰影瞬時空間效應與都市熱慣性累積效應兩種不同時間尺度物理機制的架構關鍵。

\begin{figure}[htbp]
    \centering
    \includegraphics[width=\textwidth]{../urban-thermal-gnn/04_training/figures/fig_overview_pistgnn_white.png}
    \caption[PI-ST-GNN 模型架構總覽]{PI-ST-GNN 模型架構總覽}
    \label{fig:overview_pistgnn}
\end{figure}

上圖呈現異質圖建構、編碼、解碼三階段，及 RGCN 擴散＋LSTM 遞迴之逐時步展開細節。


\subsection{關係圖卷積網路（RGCN）}
\label{subsec:rgcn}

本研究的核心空間建模模組為\textbf{關係圖卷積網路}（Relational Graph Convolutional Network, RGCN） \citep{wu2023data, shan2025physics}。不同於同質圖卷積（GCN）以單一共享權重聚合所有鄰域，RGCN 為每種關係類型 $r \in \mathcal{R}$ 維護獨立的可學習投影矩陣 $\mathbf{W}_r$，得以在訊息傳遞過程中區分不同物理因果關係之訊號貢獻。本研究之關係型聚合公式與殘差＋LayerNorm 更新規則，採用並延伸 \citet{xin2025urbangraph} 於 UrbanGraph 架構中提出之關係圖卷積形式，於此基礎上進一步將節點集合區分為建物與空氣兩種異質節點類型（見第~\ref{subsec:nodes}節），並以獨立編碼器分別處理其靜態／動態特徵。

對於節點 $i$，第 $l$ 層 RGCN 塊（RGCNBlock）的更新規則定義為：

\begin{equation}
    \mathbf{h}_i^{(l+1)} = \sigma\!\left( \mathbf{W}_{\text{self}} \mathbf{h}_i^{(l)}
    + \sum_{r \in \mathcal{R}} \frac{1}{\left|\mathcal{N}_r(i)\right|}
    \sum_{j \in \mathcal{N}_r(i)} \mathbf{W}_r \mathbf{h}_j^{(l)} \right) + \mathbf{h}_i^{(l)}
    \label{eq:rgcn}
\end{equation}

其中 $\mathbf{W}_{\text{self}}$ 為自迴路權重矩陣、$\mathbf{W}_r$ 為關係 $r$ 專屬投影矩陣（本研究 $|\mathcal{R}| = 5$）、$\mathcal{N}_r(i)$ 為節點 $i$ 在關係 $r$ 下的鄰居集合、$\sigma(\cdot)$ 為 PReLU 激活函數，末項為殘差連接並於激活後接 LayerNorm 與 Dropout（$p=0.1$）。模型採用三層堆疊（$L=3$）的 RGCNBlock，允許訊息在三跳鄰域內擴散，有效涵蓋 $80\times80$ 公尺場地內建築陰影的最大投射距離；各關係矩陣之完整維度設定見附錄~\ref{appx:rgcn_detail}。此一訊息聚合機制之圖示見圖~\ref{fig:rgcn_detail}。

\begin{figure}[htbp]
    \centering
    \includegraphics[width=0.85\textwidth]{../urban-thermal-gnn/04_training/figures/fig_rgcn_detail.png}
    \caption[RGCN 訊息聚合機制示意圖]{RGCN 模組訊息聚合機制示意圖}
    \label{fig:rgcn_detail}
\end{figure}

上圖中央目標節點 $i$ 接收來自五類關係型鄰居節點的加權訊息，每類關係使用獨立投影矩陣 $\mathbf{W}_r$ 轉換後進行均值聚合，最終與自迴路項疊加並通過殘差＋LayerNorm 輸出；三層堆疊使訊息傳遞最遠可達三跳距離之鄰域。


\subsection{全域情境融合與 LSTM 時序建模}
\label{subsec:lstm}

場景整體氣候狀態（如晴天 GHI 強度、背景風速）是影響全場域 UTCI 分佈之關鍵驅動因素，但此資訊難以從局部節點特徵中推導，因此需要一個獨立於局部訊息傳遞之外的全域情境通道。全域情境編碼器對每個時步 $t$ 提取 8 維全場域統計特徵，以三層 MLP 映射至 192 維情境向量：

\begin{equation}
    \mathbf{c}_t = \text{MLP}_{\text{ctx}}\!\left( \left[\mu_{T_a}^t,\ \mu_{V_a}^t,\ \mu_{\text{RH}}^t,\ \text{GHI}^t,\ \Delta\mu_{T_a}^t,\ \Delta\mu_{V_a}^t,\ \Delta\text{GHI}^t,\ \text{sol\_alt}^t \right] \right)
    \label{eq:ctx}
\end{equation}

對每個空氣節點，Fusion MLP 融合 RGCN 輸出之節點特徵（256維）與全域情境向量（192維），投影至 256 維融合表示 $\mathbf{f}_{i,t}$，再送入共享 LSTM 捕捉 UTCI 之日間熱積累動態：

\begin{equation}
    \mathbf{o}_{i,t},\ \mathbf{h}_{i,t}^{\text{LSTM}} = \text{LSTM}\!\left( \mathbf{f}_{i,t},\ \mathbf{h}_{i,t-1}^{\text{LSTM}} \right)
    \label{eq:lstm}
\end{equation}

LSTM 之初始隱藏狀態由 $t=0$ 時刻的 RGCN 輸出投影得出而非零向量，使模型能以場景幾何資訊初始化時序動態的起始狀態，避免冷啟動問題；此一暖啟動機制沿用 \citet{xin2025urbangraph} 之設計。最終，輸出 MLP 將 LSTM 輸出映射至 UTCI 預測值 $\hat{u}_{i,t}$（z-score 正規化空間中預測，推理後去正規化還原至實際溫度單位 °C）。完整超參數設定見附錄~\ref{appx:model_hparams}。


\section{物理引導損失函數}
\label{sec:physics_loss}

純資料驅動之損失函數僅要求預測值逼近模擬標籤，並不保證預測結果服從基本熱力學原則——模型仍可能在資料稀疏區域輸出陰影節點溫度高於日照節點的物理不可能結果。本研究因此於資料損失之外，引入三項\textbf{物理資訊損失函數}（Physics-Informed Loss）作為訓練過程中的軟約束監督機制 \citep{shan2025physics}，總損失函數定義為：

\begin{equation}
    \mathcal{L}_{\text{total}} = \mathcal{L}_{\text{data}} + \alpha_1 \mathcal{L}_{\text{rad}} + \alpha_2 \mathcal{L}_{\text{temp}} + \alpha_3 \mathcal{L}_{\text{wind}}
    \label{eq:total_loss}
\end{equation}

其中 $\mathcal{L}_{\text{data}}$ 為資料驅動項（與物理式模擬標籤之 MSE 損失），$\alpha_1, \alpha_2, \alpha_3$ 為物理損失權重係數。三項物理約束各自對應之空間、時間、風場作用機制與懲罰觸發條件，統整於圖~\ref{fig:physics_loss}；各懲罰項之完整符號定義與數值設定說明如下三小節，其嚴謹解析式推導則置於附錄~\ref{appx:physics_loss_derivation}。

\begin{figure}[htbp]
    \centering
    \includegraphics[width=0.9\textwidth]{../urban-thermal-gnn/04_training/figures/fig_physics_loss.png}
    \caption[物理引導損失函數示意圖]{物理引導損失函數示意圖}
    \label{fig:physics_loss}
\end{figure}

上圖三項約束分別作用於不同物理維度：輻射空間一致性約束（左）比較日照／陰影節點群組之預測值相對關係；都市熱慣性時序平滑約束（中）限制相鄰時步 UTCI 之最大變化率；風場阻滯尾流約束（右）抑制高層建築背風側之異常熱累積。灰色區域代表各約束實際產生懲罰梯度之觸發範圍，其餘區域損失為零，僅在物理不可能之預測情形下介入訓練。


\subsection{輻射空間一致性約束（$\mathcal{L}_{\text{rad}}$）}
\label{subsec:radiation_penalty}

建築陰影投射是影響行人 UTCI 最關鍵的幾何因子之一：日照節點之熱壓力應顯著高於陰影節點，而非兩者相近或反轉 \citep{coccolo2018thermal}。當太陽高度角 $\theta_{\text{sol}} > 10^\circ$ 時，依 SVF 與即時陰影旗標將節點分類為日照節點（$\text{SVF}_i > 0.5$ 且未受遮蔽）與陰影節點（$\text{SVF}_i < 0.5$ 且受遮蔽），懲罰項要求場域陰影節點均值與日照節點均值之間至少維持最小間隔 $m_{\text{rad}}$，對未達此間隔（含反轉）之部分施以平方懲罰：

\begin{equation}
    \mathcal{L}_{\text{rad}} =
    \left[ \text{ReLU}\!\left( \bar{u}_{\text{shade}} - \bar{u}_{\text{sun}} + m_{\text{rad}} \right) \right]^2
    \label{eq:lrad}
\end{equation}

其中 $\bar{u}_{\text{shade}}$、$\bar{u}_{\text{sun}}$ 分別為當前時步場域陰影節點與日照節點之預測均值，最小間隔 $m_{\text{rad}} = 0.5$（對應正規化 UTCI 空間中約半個標準差）；式~\eqref{eq:lrad}於每個滿足 $\theta_{\text{sol}}>10^\circ$ 之時步各計算一次，訓練時取其對有效時步之平均。此設計要求日照節點至少比陰影節點高出 $m_{\text{rad}}$，而非僅禁止陰影反超日照——理由是完全日照與完全遮蔭之間應存在物理上有意義的溫差，僅約束「不反轉」仍可能容許模型學習出兩者幾乎相等的退化解。


\subsection{都市熱慣性時序平滑約束（$\mathcal{L}_{\text{temp}}$）}
\label{subsec:temporal_penalty}

都市物質（混凝土、瀝青鋪面）具有顯著的熱儲存效應，使 UTCI 在相鄰時步間的變化率受到自然物理限制——熱壓力不應在一小時內劇烈跳躍。本研究將允許之最大 UTCI 變化率設定為 $5^\circ\text{C/hr}$（正規化空間中 $\delta_{\text{temp}} = 0.625$，其換算依據見附錄~\ref{appx:physics_loss_derivation}）：

\begin{equation}
    \mathcal{L}_{\text{temp}} = \text{mean}_{i,t}\!\left[ \text{ReLU}\!\left( \left| \hat{u}_{i,t+1} - \hat{u}_{i,t} \right| - \delta_{\text{temp}} \right) \right]
    \label{eq:ltemp}
\end{equation}

式中對超出容許變化率之部分採\textbf{線性}鉸鏈懲罰（而非平方），使梯度大小不隨違反程度二次放大，讓訓練初期時序尚不穩定時，此一約束不至於因梯度過大而主導總損失。此懲罰項能有效防止模型在學習日出／日落高輻射梯度時產生非物理的時間振盪，提升整體逐時預測軌跡的熱力學自洽性 \citep{oskarsson2023modeling}。


\subsection{風場阻滯尾流約束（$\mathcal{L}_{\text{wind}}$）}
\label{subsec:wind_penalty}

高層建築（$B_h > 20\text{m}$，即正規化特徵 $B_h/50 > 0.4$）在背風側產生顯著的渦旋死角（Leeward Wake），使局部熱能累積形成「熱死角」現象 \citep{oke1987boundary}。對於位於高層建築鄰近範圍之空氣節點集合 $\mathcal{B}_{\text{tall}}$，風場阻滯懲罰定義為：

\begin{equation}
    \mathcal{L}_{\text{wind}} = \text{mean}_{i \in \mathcal{B}_{\text{tall}}}
    \left[ \text{ReLU}\!\left( \hat{u}_i - \mu_{\hat{u}} - 1.5\,\sigma_{\hat{u}} \right) \right]
    \label{eq:lwind}
\end{equation}

其中 $\mu_{\hat{u}}$、$\sigma_{\hat{u}}$ 分別為當前時步全場域預測 UTCI（所有節點，非僅 $\mathcal{B}_{\text{tall}}$）之均值與標準差，懲罰同樣採線性鉸鏈形式。閾值 $\mu_{\hat{u}} + 1.5\sigma_{\hat{u}}$ 之選定依循統計離群偵測之常見慣例——在常態分布假設下，超出均值 $1.5$ 個標準差約對應最熱的 6.7\% 尾端節點，將「熱死角」操作型定義為場域中溫度顯著偏離常態分布中心之高層建築鄰近節點。此約束強制模型學習高層建築背風側之熱累積模式，使設計師能可視化熱死角位置，進而最佳化建築間距與通風廊道配置。


\section{模組五：Grasshopper 介面整合與多目標最佳化——熱舒適人行動線設計}
\label{sec:nsga2}

\subsection{從最短路徑到熱舒適動線評估}
\label{subsec:from_astar_to_walkway}

早期構想以 $A^*$ 最短路徑演算法將 UTCI 熱壓力場轉譯為單一起訖點間之熱迴避路徑：以節點間歐氏距離與熱壓力懲罰之加權組合作為邊成本，求解一條在行走效率與熱舒適間取捨之路徑。然而，都市開放空間設計決策的對象通常並非單一使用者的單次行程，而是整體人行動線網絡在一整日各時段的熱舒適表現；單點對單點的最短路徑求解難以回答「這個街廓的人行動線系統整體而言是否宜行」此一設計層級問題。本研究因此將原先的路徑搜尋機制，重新定位為\textbf{熱舒適人行動線設計}（Thermal Comfort-Based Walkway Design）評估指標，直接整合進 NSGA-II 多目標最佳化之適應度函數，而非作為獨立之單次路徑查詢工具。

具體而言，設計者於 Grasshopper 中定義之人行動線網絡幾何（步行軸線、廣場動線、建物出入口連通路徑）作為固定拓樸輸入，GNN 代理模型即時預測之全場域 UTCI 熱壓力場則沿此網絡逐段取樣，計算全動線之\textbf{時空平均熱暴露指標}：

\begin{equation}
    \bar{\Phi}_{\text{walkway}} = \frac{1}{T}\sum_{t=1}^{T} \frac{1}{|E_{\text{walk}}|}
    \sum_{(i,j) \in E_{\text{walk}}} D_{ij}\cdot \Phi\!\left(\text{UTCI}_{j,t}\right)
    \label{eq:walkway_exposure}
\end{equation}

其中 $E_{\text{walk}}$ 為人行動線網絡之邊集合、$D_{ij}$ 為邊之實際步行距離、$\Phi(\cdot)$ 為正規化熱壓力懲罰函數，其定義與完整推導見附錄~\ref{appx:walkway_derivation}。此一時空平均指標取代原先僅評估單一路徑之作法，使評估對象與 NSGA-II 之全域最佳化目標在尺度上一致。


\subsection{Grasshopper 介面與非同步通訊架構}
\label{subsec:gh_integration}

系統採用\textbf{非同步用戶端-伺服器架構}，讓建築師能在熟悉的設計環境中獲得 AI 的即時反饋：在 Grasshopper 中開發之專用套件（\texttt{UTCIPredictor}、\texttt{UTCIOptimizer}）負責將 Rhino 3D 幾何物件序列化為異質圖相容張量；後端透過 IETF RFC 6455 \textbf{WebSocket} 通訊協定 \citep{rfc6455} 與 FastAPI 通訊，載入預先訓練完成的 PI-ST-GNN 模型權重回傳預測結果與最佳化進度；設計師於 Rhino 中移動建築時，系統更新場地 1.0m 間距之高解析度熱舒適度網格，使原本須另行呼叫模擬引擎之熱舒適評估得以整合入設計操作流程 \citep{sadeghipour2013ladybug}。

圖~\ref{fig:gh_websocket} 以時序圖具體呈現此一通訊架構之兩種訊息流：(a) 單次幾何更新之同步請求-回應循環，直接對應每次推論之低延遲特性；(b) NSGA-II 最佳化工作之非阻塞式執行——最佳化迴圈於獨立執行緒（\texttt{ThreadPoolExecutor}）中運行，避免長時間運算阻塞 WebSocket 事件迴圈，並透過 \texttt{asyncio.Queue} 將逐代進度（世代數、當前最佳解、可行解數量）以非同步訊息即時串流回 Grasshopper 前端，使建築師無需等待完整最佳化流程結束即可觀察收斂趨勢。

\begin{figure}[htbp]
    \centering
    \includegraphics[width=\textwidth]{../urban-thermal-gnn/04_training/figures/fig_gh_websocket_sequence.png}
    \caption[非同步通訊架構時序圖]{Grasshopper--FastAPI／WebSocket 非同步通訊架構時序圖}
    \label{fig:gh_websocket}
\end{figure}

上圖實線為同步 WebSocket 訊息，虛線為經 \texttt{asyncio.Queue} 橋接執行緒之非同步進度回呼。


\subsection{染色體編碼機制與多目標函數}
\label{subsec:design_variables}

最佳化引擎將都市配置參數化為連續染色體基因組，所有變數於運算前歸一至 $[0,1]$ 區間：每棟建築由質心座標 $(c_x, c_y)$、面寬 $w$、進深 $d$、平面旋轉角度與樓層數共六項參數定義；每棵喬木由座標 $(x,y)$、樹冠半徑與高度四項參數組成，直接決定遮蔭效果與綠化率貢獻。

系統透過 NSGA-II \citep{deb2002fast} 同時處理三個相互競爭的適應度指標：

\begin{enumerate}
    \item \textbf{$f_1$（全域 UTCI 極小化）：}最小化全場地所有開放空間空氣節點在 11 時段內的時空平均熱壓力。
    \item \textbf{$f_2$（人行動線熱暴露極小化）：}最小化式~\eqref{eq:walkway_exposure}定義之人行動線時空平均熱暴露指標 $\bar{\Phi}_{\text{walkway}}$。
    \item \textbf{$f_3$（綠化率極大化）：}最大化場地綠化率，兼顧生態價值。
\end{enumerate}

法規合規性透過 5 維違規向量（FAR、BCR、法定縮退、場地容納性、建築碰撞）加以監控，採「可行性優先」選擇啟發式：符合規範之可行解絕對優先於違規解排序。最終輸出之帕累托前緣代表一組最優設計集，讓建築師在整體熱舒適性能、人行動線品質與綠化率之間進行數據驅動的決策。

本研究之最佳化引擎（\texttt{07\_optimization/nsga2\_engine.py}）已完整實作 $f_1$（全域 UTCI）與 $f_3$（綠化率）兩項目標之適應度函數，並以此二維目標空間驗證第~\ref{sec:pareto_analysis}節之收斂表現；$f_2$（人行動線熱暴露）之數學定義與推導已於式~\eqref{eq:walkway_exposure}及附錄~\ref{appx:walkway_derivation}完整給出，其整合入適應度評估器（\texttt{FitnessEvaluator}）之工程實作——即將染色體所定義之人行動線網絡幾何與 GNN 預測場逐段取樣、回傳第三項目標值——列為第~\ref{sec:future_work}節之近期後續工作，目前尚未併入單次最佳化執行流程，特此說明，避免與第~\ref{sec:pareto_analysis}節實際呈現之收斂結果所涵蓋目標數產生混淆。


\subsection{與 Wallacei 都市型態最佳化案例之比較}
\label{subsec:wallacei_comparison}

Grasshopper 生態系中，Wallacei 是目前應用最廣泛之進階 NSGA-II 視覺化最佳化外掛，已見於多項都市型態多目標最佳化實務研究 \citep{wortmann2017opossum, daglier2025developing}。既有以 Wallacei 為核心之都市形態最佳化案例，其適應度函數之環境績效評估多直接呼叫 Ladybug Tools／EnergyPlus 進行逐代個體之全氣候模擬，因此在演算法收斂所需的數百代、每代數百個體之規模下，單次最佳化任務之總運算時間常達數週甚至數月，實務上多僅能於方案設計完成後進行事後驗證，而非設計初期之即時決策引導 \citep{aman2023ai, daglier2025developing}。

本研究之最佳化引擎與 Wallacei 案例共用相同的 NSGA-II 核心演算法與 Grasshopper 介面整合模式，惟以第~\ref{sec:gnn_model}節訓練完成之 PI-ST-GNN 取代 Wallacei 案例中直接呼叫之高擬真度模擬引擎作為適應度函數之環境績效來源。此一取代之目的並非單純追求運算速度提升，而在於使多目標最佳化之搜尋迴圈得以在設計對話的時間尺度內完成，讓建築師於方案概念階段即可即時比較不同人行動線配置與量體組合之熱舒適表現，而非僅能於方案定案後進行單次驗證。兩者於相同法規約束與設計變數空間下之收斂表現比較，見第~\ref{chapter:results}章。
